\documentclass[aspectratio=169]{beamer}%[handout]
\input{Macros}

\def\titulo{Título}
\def\subtitulo{Subtítulo}
\def\autora{Autor}
\def\correoa{correo@ulagos.cl}
\def\autorb{Autor}
\def\correob{correo@ulagos.cl}
\def\autorc{Autor}
\def\correoc{correo@ulagos.cl}
\def\asignatura{Asignatura}
\def\campus{Campus}
\def\carrera{Ingeniería Civil en Informática}

\lstset{language=SQL}

%Unico AUTOR
%\author[\autora]{\hspace{5cm}\raggedright\autora\\\hspace{5cm}\href{mailto:\correoa}{\correoa}}
   \author[NOMBRE AUTOR]{
	\href{mailto:\correoa}{\autora}\\
	\href{mailto:\correob}{\autorb}\\
	\href{mailto:\correoc}{\autorc}}
\institute[ULAGOS]{\hspace{5cm}\raggedright{\carrera}}
%%%%%%%%%%%%%FIN PORTADA


%%%%%Carpeta predeterminada fotos


\begin{document}


%Pagina de Portada
\begin {frame} [plain]

\titlepage
\end {frame}
\setbeamertemplate{background}{}
%%%%%%%%%%%%%%%%ÍNDICES
\section[Contenido]{}
\frame{
  \frametitle{\textbf{Contenido}}
\setcounter{tocdepth}{2}%1: solo titulo principal, 2: titulo y subtitulo, 3....
\scriptsize
\tableofcontents[]
}%Generacion de Indice por capitulo
\AtBeginSection[]{
\begin{frame}
\frametitle{\textbf{Contenido}}
\scriptsize
\tableofcontents[currentsection]
\end{frame}}

\AtBeginSubsection[]{
\begin{frame}
\frametitle{\textbf{Contenido}}  
\scriptsize
\tableofcontents[currentsection,currentsubsection]
\end{frame}
}
%%%%%%%%%%%%%%%%FIN ÍNDICES







%%%%%%%%%%%%%%%%%%%%%%%%%%%%%%%%%%%%%
%%%%%%%%%%%%%%%%%INICIO PRESENTACION
\section{Introducción}
\section{Marco Teórico}
\section{Estado del Arte}
\section{Justificación}
\section{Objetivos}
\section{Desarrollo}
\subsection{Metodología de Desarrollo}
\subsection{Modelo de Datos}
\subsection{Requisitos}
\begin{frame}[fragile]
	\frametitle{Colores y Bloques}
	\justifying
	Lo siguiente son ejemplos de texto en colores: \rojo{Hola}, \verde{Hola}, \azul{hola}, \naranjo{hola}
	
	\begin{block}{Titulo}\justifying
		Contenido
	\end{block}
	
	
	\begin{exampleblock}{Titulo ejemplo}\justifying
		Contenido
	\end{exampleblock}
	
	
	\begin{alertblock}{Titulo alerta}\justifying
		Contenido
	\end{alertblock}
	
	\begin{sintaxis}\justifying
	\end{sintaxis}
	
	\begin{resultado}\justifying
	\end{resultado}
	
\end{frame}


\begin{frame}[fragile]
	\frametitle{Items, Descripciones y Enumeraciones}
	\justifying
	\transdissolve
	\begin{itemize}\justifying
		\item Item sin números
		\begin{itemize}\justifying
			\item nivel 2
			\begin{itemize}\justifying
				\item nivel 3
			\end{itemize}
		\end{itemize}
	\end{itemize}
	
	\begin{enumerate}\justifying
		\item Item Numerado
		\begin{enumerate}\justifying
			\item Nivel 2
			\begin{enumerate}\justifying
				\item Nivel 3
			\end{enumerate}
		\end{enumerate}
	\end{enumerate}
	
	\begin{description}\justifying
		\item[Descripción] Texto descrito
		\begin{description}\justifying
			\item[Nivel 2] Texto
		\end{description}
		
	\end{description}
\end{frame}




\subsection{Códigos}
\begin{frame}[fragile]
	\justifying
	\frametitle{Código}
	\lstset{language=C}
	\begin{lstlisting}
		#include<stdio.h>
		int main(){
			printf("Hola Mundo1");
			return 0;
		}
	\end{lstlisting}
	\begin{block}{SQL}%\vspace{-0.3cm}
		\lstset{language=SQL}
		\begin{lstlisting}
			Select *
			from navegantes n
			where categoria=7;
		\end{lstlisting}%\vspace{-0.3cm}
	\end{block}
	
\end{frame}

\begin{frame}[fragile]
	\frametitle{Código multicolumna}
	\justifying
	\lstset{language=C}
	\begin{minipage}[t]{0.49\textwidth}
		\begin{lstlisting}
			#include<stdio.h>
			int main(){
				printf("Hola Mundo2");
				return 0;
			}
		\end{lstlisting}
	\end{minipage} \hfill
	\begin{minipage}[t]{0.49\textwidth} 
		\lstset{firstnumber=6} %cambiar numero de inicio
		\begin{lstlisting}
			#include<stdio.h>
			int main(){
				printf("Hola Mundo3");
				return 0;
			}
		\end{lstlisting}
	\end{minipage}
\end{frame}
\subsection{Matemática}
\begin{frame}[fragile]
	\frametitle{Ejemplo de Función Matemática}
	\justifying
	\begin{displaymath}
		C_L=\frac{(S_{22}-\delta S_{11}^*)^*}{|\varPi S_{22}|^2=-|\pi|^2}
	\end{displaymath}
	
	\begin{displaymath}
		R_S=\frac{\sqrt{1-g_s}\cdot (1-|S_{11}|^2)}{1-(1-g_s)\cdot|S_{11}|^2}
	\end{displaymath}
	
\end{frame}



\begin{frame}[fragile]
	\frametitle{Tabla}
	\justifying
	\begin{center}
		\begin{tabular}{|c|c|c|c|c|c|}\hline
			\textbf{S}&\textbf{SCT} &\textbf{Asignatura}&\multicolumn{2}{|c|}{\textbf{Total Horas}}&\textbf{Previatura} \\\cline{4-5}
			&&&\textbf{TP}&\textbf{TA}&\\\hline
			\multirow{5}{*}&a&b&c&d&r\\\hline
		\end{tabular}
	\end{center}
\end{frame}

\section{Conclusión}
\section{Trabajo Futuro}


%%%%%%%%%%%%%%%%%%%%%%%%%%%%%%
%%%%%%%%%%%%%%%%%FIN DOCUMENTO
\end{document}
